\section{Method}
\label{sec:method}

\subsection{Amortized objective}

Let $L(w;o,\sigma)$ be the solver's objective at one pixel. We train one network $f_\theta$ over the pixels of the training set,
\begin{equation}
  \theta^\star=\argmin_\theta\ \sum_{p} L\bigl(f_\theta(o_p);\,o_p,\sigma_p\bigr).
  \label{eq:amortized}
\end{equation}
No solver DEM enters \cref{eq:amortized}. The network sees only the six intensities of one pixel, neither its neighbors nor $\sigma$, which the loss alone uses (\cref{fig:pipeline}).

\begin{figure}[t]
\centering
\begin{tikzpicture}[
  font=\footnotesize,
  box/.style={draw=black!60, rounded corners=1.5pt, minimum height=5.6mm, inner sep=2.5pt, align=center},
  loss/.style={box, fill=black!6},
  arr/.style={-{Stealth[length=4pt]}, black!70, line width=0.4pt}]
  \node[box] (o) {$o$\\[-1pt]{\scriptsize 6 intensities}};
  \node[box, right=4.5mm of o] (f) {$\sqrt{\cdot}$ + Fourier\\[-1pt]{\scriptsize MLP, SiLU}};
  \node[box, right=4.5mm of f] (w) {$w\ge0$\\[-1pt]{\scriptsize 54 weights}};
  \node[box, right=4.5mm of w] (x) {DEM $Bw$\\[-1pt]{\scriptsize 18 bins}};
  \node[box, below=4.5mm of w] (r) {$Dw$\\[-1pt]{\scriptsize predicted intensities}};
  \node[loss, left=6mm of r] (l) {$\sum_k w_k+\mu\sum_c (v_c/h_c)^2$\\[-1pt]{\scriptsize BP loss, no solver output}};
  \draw[arr] (o) -- (f);
  \draw[arr] (f) -- (w);
  \draw[arr] (w) -- (x);
  \draw[arr] (w) -- (r);
  \draw[arr] (r) -- (l);
  \draw[arr] (o.south) -- node[right, inner sep=1.5pt] {\scriptsize $o,\sigma$} (o.south |- l.north);
\end{tikzpicture}
\caption{Training. The network maps one pixel's six intensities to the solver's 54 weights. The loss compares the predicted intensities with the observation and its noise band and adds the total weight.}
\label{fig:pipeline}
\end{figure}

We take the square root of each intensity and lift it to sines and cosines at 12 fixed frequencies~\cite{tancik2020fourier,mildenhall2020nerf}, giving 150 inputs together with the roots themselves. Four hidden layers with SiLU activations~\cite{elfwing2018silu} feed a linear layer and a softplus that keeps the weights non-negative. The final BP model is 680 wide (1.5M parameters) and the Elastic Net model 336 wide (409k). The network outputs the solvers' own variable, the 54 weights, instead of the 18 DEM bins, so it searches the same space as the solver and BP's sparsity term means the same thing for both.

\subsection{Losses}

For Elastic Net, $L$ is \cref{eq:enet} itself. For BP we need a differentiable objective for any output, including outputs outside the band, and the hard constraint becomes a squared hinge:
\begin{equation}
  L_{\mathrm{BP}}(w)=\sum_{k} w_k+\mu\sum_{c}\Bigl(\frac{v_c(w)}{h_c}\Bigr)^2,
  \label{eq:loss-bp}
\end{equation}
where $v_c(w)=\max\bigl(0,\,|(Dw)_c-o_c|-h_c\bigr)$ is the distance outside the band in DN/s. Inside the band the penalty vanishes, so no fit term remains and the sparsity term breaks ties as it does in BP. Dividing by $h_c$ makes one half-width outside cost the same in every channel. Both losses are convex in $w$ for each pixel.

\subsection{What the penalty weight does}
\label{sec:law}

$L_{\mathrm{BP}}$ is the quadratic exterior penalty of the LP~\eqref{eq:bp}~\cite{courant1943variational,fiacco1968nonlinear,nocedal2006numerical}. The textbook result is that its minimizer approaches the LP's solution as $\mu\to\infty$, and Mangasarian~\cite{mangasarian2004newton} showed that a large enough finite weight already gives an exact dual solution. Training needs a more specific statement: where does a pixel settle at a given $\mu$, and when does that answer stop being accurate?

Let $y^\star$ be the LP's dual solution, positive where a lower band edge binds and negative where an upper edge binds, and let $\lambda_c=|y^\star_c|$, the total weight saved per DN/s by which that edge is relaxed. Let $d_c=v_c/h_c$ be the violation in half-widths.

\begin{proposition}
\label{prop:law}
Suppose the LP~\eqref{eq:bp} is feasible with unique primal and dual solutions $w^\star$ and $y^\star$, with support $S$ and binding channels $A$, and let $q=D_{A,S}^{-1}(y^\star_A\circ h_A^2)$. There is a finite $\mu_0$ such that for every $\mu\ge\mu_0$ a minimizer of \cref{eq:loss-bp} is zero off $S$, equals
\begin{equation}
  w_{\mu,S}=w^\star_S-\frac{q}{2\mu}
  \label{eq:path}
\end{equation}
on $S$, and has in every channel
\begin{equation}
  d_c=\frac{\lambda_c h_c}{2\mu}=\frac{\lambda_c\,t\sigma_c}{2\mu}.
  \label{eq:law}
\end{equation}
Moreover $\mu_0$ is the largest $\mu$ at which \cref{eq:path} stops being valid: an element reaches zero, at $\mu=q_k/(2w^\star_k)$, or a free channel reaches a band edge.
\end{proposition}

\Cref{app:proof} gives the proof, but the law itself is a balance of gradients. Stepping $v$ outside the band saves $\lambda v$ in total weight at a constant rate, while the penalty $\mu(v/h)^2$ has slope $2\mu v/h^2$, and the minimizer sits where the two slopes match. A pixel therefore settles further out when light in that channel is expensive, when the reading is noisy, and when $\mu$ is small.

\paragraph{Prices.}
When a single edge binds, the LP buys light from the most efficient basis element, and $\lambda_c=1/\max_k D_{ck}$, the single-edge price of \cref{tab:threshold}. A unit of weight buys at most 0.68\,DN/s of 94\Ang{} but about 290\,DN/s of 171\Ang, so 94\Ang{} is over 400 times more expensive. Real pixels usually bind several edges, and their prices can differ from this value.

\paragraph{From a tolerance to $\mu$.}
Inverting \cref{eq:law}, a tolerance of $\varepsilon$ half-widths needs $\mu\ge\lambda_c h_c/(2\varepsilon)$. With the single-edge price this takes only $R$, $B$ and $\sigma$, with no solver and no training. Take a typical flare core with a 94\Ang{} reading of 164\,DN/s and $\sigma=11.5$, so $h=16.1$ and the band runs from 148 to 180. \Cref{eq:law} gives $d=11.9/\mu$ half-widths, and $\varepsilon=0.1$ needs $\mu\approx120$. Because $\lambda_c$ and $\sigma_c$ vary, no single $\mu$ suits every pixel, so the law does not pick a value but tells us what each value buys. At feasible pixels a larger $\mu$ is always closer to the LP, which leaves training as the upper limit on $\mu$.

\paragraph{Dropping plasma.}
Below $\mu_0$ the law breaks because an element leaves the solution or a free channel reaches an edge, and \cref{eq:path} says when. A one-element model shows how this scales. Take a channel whose signal can only be bought from an element of yield $r=1/\lambda$, with band floor $\ell=o-h$. Minimizing $w+\mu\bigl((\ell-rw)/h\bigr)^2$ over $w\ge 0$ gives
\begin{equation}
  r w^\circ=\max\Bigl(0,\ \ell-\frac{\lambda h^2}{2\mu}\Bigr),
  \label{eq:clip}
\end{equation}
and all of that plasma is dropped once
\begin{equation}
  \mu\le\mudrop=\frac{\lambda h^2}{2\ell}.
  \label{eq:drop}
\end{equation}
For a shot-noise-limited reading, $\sigma^2=g\,o$ with $g$ fixed by the channel and exposure time. Then for $o\gg h$,
\begin{equation}
  \mudrop\approx\frac{\lambda\,t^2 g}{2},
  \label{eq:drop-shot}
\end{equation}
which does not depend on brightness: a brighter pixel needs proportionally more plasma, but the penalty for leaving it out, $(o/h)^2=o/(t^2 g)$, grows at the same rate. At typical noise only 94\Ang{} exceeds 1 (\cref{tab:threshold}), and at the flare core above $\mudrop=1.3$. At $\mu=1$ the first unit of plasma costs 1 and saves only 0.77 of penalty, so the minimum buys none for this channel, whereas at $\mu=100$ it buys until the reading is 0.12 half-widths, about 2\,DN/s, short of the band.

\begin{table}[t]
\centering
\small
\caption{Per-channel single-edge price in weight units per DN/s, noise ratio $g=\sigma^2/o$ (median over bright, high-signal validation pixels), and the resulting threshold \eqref{eq:drop-shot}. Short flare exposures raise $g$ and the threshold in some channels.}
\label{tab:threshold}
\begin{tabular}{@{}lcccccc@{}}
\toprule
Channel (\AA) & 94 & 131 & 171 & 193 & 211 & 335\\
\midrule
Price & 1.48 & 0.10 & 0.0034 & 0.0070 & 0.026 & 0.60\\
$g$ & 0.81 & 0.81 & 0.76 & 0.58 & 0.38 & 0.73\\
\mudrop & \textbf{1.18} & 0.08 & 0.003 & 0.004 & 0.009 & 0.43\\
\bottomrule
\end{tabular}
\end{table}

\paragraph{Infeasible pixels.}
If no $w$ fits every band, the violations cannot vanish and \cref{eq:law} does not apply. A larger $\mu$ then makes the penalty dominate the loss, and the minimizer tends to the closest fit at any weight, whereas the solver returns the smallest weight inside a widened band.

\subsection{Separating loss failures from training failures}
\label{sec:gap}

Because $L_{\mathrm{BP}}$ is convex per pixel, we can find its minimizer $w^\circ(o,\sigma)$ numerically without a network, using L-BFGS-B, and no network can score better on that pixel. This gives two tests. The \emph{amortization gap}~\cite{cremer2018inference} compares the network's loss with that of $w^\circ$: if the gap is small and the result is poor, the loss itself is at fault and more training cannot help. The \emph{screen} scores $w^\circ$ against the solver's DEM for each candidate $\mu$, evaluating a loss on held-out pixels without training a model.

\subsection{Clamped channels}
\label{sec:clamp}

PSF deconvolution can produce negative intensities, and these are then set to zero. We call a channel clamped when it reads at most $10^{-3}$\,DN/s, and 40\% of valid pixels in our data have such a channel, nearly all of them faint. A clamped reading is an artifact, so a band around it asks for the wrong thing. Dropping the whole pixel, as we first did, removes every input containing a zero from the training set. \emph{Channel masking} instead keeps the pixel and removes only that channel's penalty: the other channels keep their bands, and the network still reads the zero as an input.

\subsection{Training}

Each epoch draws 512 fresh pixels from each of 58,688 blocks of $128\times128$ pixels, about 30M pixels, in batches of 16 blocks. We use Adam~\cite{kingma2015adam} at a peak learning rate of $3\times10^{-4}$ with no weight decay, a linear warm-up (500 steps for BP, 3,000 for Elastic Net), cosine decay and gradient-norm clipping at 1. The final BP model trains for 120 epochs with $\mu=100$ and channel masking, in 6.6 hours on one GPU. The Elastic Net model trains for 40 epochs without clamped pixels. \Cref{app:engineering} describes a silent collapse at the start of training and a sampling bug, with their fixes.
